\section{Introduction}\label{sec:introduction}

Let $\mathbb T=\mathbb R/(2\pi\mathbb Z)$, with normalized Lebesgue
measure $d\mu(x)=dx/(2\pi)$. For an integrable complex-valued function
$f$, write
\[
 \widehat f(k)=\int_{\mathbb T} f(x)e^{-ikx}\,d\mu(x),
 \qquad
 S_N f(x)=\sum_{k=-N}^{N}\widehat f(k)e^{ikx}.
\]
The endpoint question concerns the Orlicz class
\[
 L\log L(\mathbb T)
 =\left\{f:\int_{\mathbb T}|f(x)|\log(2+|f(x)|)\,d\mu(x)<\infty\right\}.
\]

\begin{theorem}[Almost-everywhere convergence in $L\log L$]
\label{thm:main}
For every complex-valued $f\in L\log L(\mathbb T)$, there is a measurable
set $N_f$ of measure zero such that
\[
 \lim_{N\to\infty}S_Nf(x)=f(x)
 \qquad (x\in\mathbb T\setminus N_f).
\]
The limit is along the full sequence $N=0,1,2,\ldots$.
\end{theorem}

Carleson proved almost-everywhere convergence for square-integrable
functions \cite{Carleson1966}, and Hunt extended this conclusion to
$L^p$ for every $p>1$ \cite{Hunt1968}. Kolmogorov's integrable
counterexample shows that the conclusion cannot hold throughout $L^1$
\cite{Kolmogorov1923}.
Between these cases, Sj\"olin obtained convergence in
$L\log L\log\log L$ \cite{Sjolin1969}, and Soria developed larger
convergence spaces through extrapolation \cite{Soria1985,Soria1989}.
Antonov proved convergence in $L\log L\log\log\log L$
\cite{Antonov1996}. Arias de Reyna subsequently introduced the
quasi-Banach space $\mathrm{QA}$, which strictly contains both Antonov's
class and Soria's space \cite{AriasDeReyna2002}.
Carro, Masty\l{}o and Rodr\'iguez-Piazza further studied $\mathrm{QA}$
and obtained a Lorentz convergence space strictly larger than Antonov's
space \cite{CarroMastyloRodriguezPiazza2012}.
Here the iterated-log class names describe growth at infinity; the
logarithms may be regularized near zero.

Lie gave a time-frequency proof of the full-sequence maximal bound from
$\mathrm{QA}$ to weak-$L^1$, avoiding extrapolation
\cite[Corollary~1.2(5)--(7)]{Lie2013}.
Di Plinio and Fragkos proved weak-$L^p$ bounds of order $(p-1)^{-1}$
as $p\downarrow1$, together with sparse estimates and weighted
endpoint convergence results \cite{DiPlinioFragkos2025}.
The $L\log L$ conjecture is also formulated through a weak-$L^1$
maximal estimate in \cite[Definition~1.1 and Conjecture~1]{Lie2017}.
\Cref{thm:main} establishes the conjectured sufficiency of $L\log L$
for the ordinary Fourier sums along the full sequence.
Stein's Banach-space theorem also yields the corresponding Luxemburg-norm
weak-$L^1$ maximal bound; see Corollary~\ref{cor:fourier-weak-l1}.
Our proof proceeds through the good-set integral estimate described next.

\subsection{The estimate behind the convergence theorem}

The main analytic estimate concerns one nonnegative input $\rho$ of mass
one and height at most $e^d$, where $d$ is large. Its uncentered spatial
maximal density is
\[
 M\rho(x)=\sup_{\substack{I\ni x\\I\text{ a circle arc}}}
            \frac{1}{\mu(I)}\int_I\rho\,d\mu,
 \qquad 0<\mu(I)\le1.
\]
Set $h=\log\log d$ and remove the points where $M\rho>e^{4h}$. The removed set has measure at most
$Ce^{-4h}$, while on its complement we prove
\begin{equation}\label{eq:intro-good-set}
 \int_{\{M\rho\le e^{4h}\}}\sup_{N\ge0}|S_N\rho|\,d\mu
 \le Cd.
\end{equation}
The two bounds have different roles. The integral bound is summable with
the masses of the height layers of an $L\log L$ function. The exceptional
sets are summable when the layers have doubly exponential heights.
\Cref{sec:fourier} makes both statements precise and proves
\Cref{thm:main} from \Cref{eq:intro-good-set}.

Good-set integral estimates also occur in Lie's near-endpoint argument
\cite[Proposition~3.1 and (3.13)]{Lie2013}. Here the bounded-density estimate
and exceptional-set bound are quantified for the summation at the
$L\log L$ endpoint.

To obtain \Cref{eq:intro-good-set}, we first use a finite binary tree of
dyadic intervals. In unit-period coordinates, modulation multiplies the
input by $e^{2\pi iuy}$, where $u$ is an integer. At each interval, take
half the difference between the averages of this modulated input on its
two children. Along a uniformly chosen output path, this coefficient is
known before the next left-or-right choice. Its output uses the sign of
one further binary choice; a multiplier of modulus at most one may depend on the
intervening choice. Stop the path when the average of the original
input first exceeds $e^{4h}$. \Cref{prop:dyadic} bounds the integral of
the supremum over all integer modulations of the resulting sum by $Cd$, uniformly
in the tree depth. This is the finite estimate assembled by the
combinatorial part of the proof.

Translation and scale averaging of these delayed Haar kernels produce a
nonzero multiple of the Hilbert kernel, with an error controlled on the
same good set. The periodic Hilbert transform then gives the ordinary
symmetric Fourier projections. The averaging follows the dyadic
representation method of Petermichl and Hyt\"onen
\cite{Petermichl2000,Hytonen2008};
\Cref{sec:fourier} includes the kernel calculation and all limiting steps.

\subsection{The structural argument}

Fefferman's proof organized the Carleson operator by a measurable frequency
selector \cite{Fefferman1973,Fefferman1997Erratum}.
Lacey and Thiele later gave a simplified proof of its weak-$(2,2)$ bound
using phase-plane tiles and trees \cite{LaceyThiele2000}.

The finite proof must control many modulation frequencies without a loss
that grows with their number. It splits the input into pieces and, on
each piece, uses a noisy observation of the input point to identify
frequencies with nearly the same phase. The input point is sampled with
probability proportional to the density, whereas the transform is
estimated along a uniformly sampled output path. Information is counted
under the first law and converted to path estimates under the second.

Three constructions make these estimates compatible.

\begin{enumerate}[leftmargin=*,itemsep=4pt]
\item \emph{One phase region for a large set of frequencies.}
The compression lemma starts with a finite set of integer frequencies
and a uniform bound on noisy sums of independent samples from any
probability law on that set. It produces a large subset whose phase
defects, relative to one center frequency, are small on average over
one common region.
The conclusion remains valid for every probability law on the retained
subset, so later spectral weights may be chosen after the subset is
found. The proof introduces correlated samples with prescribed marginals
and controls their relative entropy together with the information in
the noisy sum. Two orders of refining their conditional laws reach the
required phase concentration with a small loss of frequencies
(\Cref{sec:entropy-foundations,sec:entropy-refinement}).

\item \emph{Two bounds for the cost of a new observation.}
For an input point $Y$, a label has the form $Z=Y+\theta$ on the circle,
where the noise density
is proportional to $e^{-A}$ and $A$ is a nonnegative linear combination
of functions $1-\cos(2\pi sy)$, $s\in\mathbb Z$. Nonnegative Fourier
coefficients and the Griffiths--Ginibre positive-correlation method
\cite{Ginibre1970} control how this density changes when a new term
is added to $A$. If too many large conditional Fourier coefficients are
separated, compression supplies such an addition. It either produces
an event unlikely under a comparison law, or enlarges the set of
frequencies whose phases vary little under the noise. The first
alternative is bounded by likelihood information, the second by a
packing estimate. Together they bound the total cost of the labels
(\Cref{sec:thermal,sec:procedure}).

\item \emph{Separate costs for information and for choosing frequencies.}
A part of the tree with fixed noise law and mass normalization is
called a live plateau; its normalization is the mean of its input piece
at entry. On each plateau a cumulative list approximates the relevant
frequencies by signed sums. At selected spatial intervals we store a
smaller list of representatives and the current mass measures divided
by that normalization; these records are the snapshots.
They are refreshed before the next output bit whenever those measures
have changed enough.

For a representative frequency, the path estimate controls a transform
whose values are functions of the label. Integrating against a label
character recovers the scalar transform at a shifted frequency, multiplied
by the corresponding Fourier coefficient of the noise density. The choice
of representative makes this coefficient close to one;
the remaining frequency discrepancy contributes a small spatial phase
error, which sums geometrically along the path.

The delayed vector estimate charges label
information with an absolute coefficient; the logarithms of the
representative counts appear only in a separate probability tail.
The distinction is needed when the local estimates are summed. The
expected costs of choosing representatives and the remaining path terms,
under uniform output sampling, are $O(d/h)$ on each live plateau. The total starting weight
of these plateaus is $O(h)$, where a starting weight is the uniform
probability of its root interval relative to the top interval, multiplied
by its mass normalization. Their product is $O(d)$. Information and
errors from input near dyadic interval boundaries are instead charged
directly across all plateaus
(\Cref{sec:capture,sec:paths,sec:assembly}).
\end{enumerate}

\subsection{Organization}

\Cref{sec:entropy-foundations,sec:entropy-refinement} prove the entropy
compression lemma. \Cref{sec:thermal} constructs and estimates the label
updates. \Cref{sec:procedure} states the finite dyadic estimate, defines its
splitting procedure, and proves the global budgets. \Cref{sec:capture} constructs frequency lists,
controls exceptional edges, and fixes the timing of snapshots.
\Cref{sec:paths} proves the scalar and vector path estimates, and
\Cref{sec:assembly} assembles the dyadic bound.
\Cref{sec:fourier} transfers that bound to Fourier sums and completes
the layer summation.
